\section{Xác định người dùng và Ma trận Phân quyền (RBAC Matrix)}
\label{sec:nguoidung_rbac}

Hệ thống ứng dụng MyHCMUT phục vụ đội ngũ nhân sự của Trường Đại học Bách khoa -- ĐHQG-HCM, được phân quyền chặt chẽ theo vai trò (Role-Based Access Control):

\begin{table}[H]
    \centering
    \renewcommand{\arraystretch}{1.3}
    \begin{tabularx}{\textwidth}{|p{3.2cm}|p{3.8cm}|X|}
    \hline
    \textbf{Tác nhân (Actor)} & \textbf{Quyền hệ thống} & \textbf{Phạm vi Chức năng trên Mobile App} \\
    \hline
    \textbf{Cán bộ / Giảng viên} & \texttt{cn:ly\_lich}, \texttt{cn:nghi\_phep}, \texttt{cn:di\_cong\_tac}, \texttt{iofficeMission:read}, \texttt{scheduleGeneral:read} & Tra cứu 11 mục lý lịch; gửi đề xuất sửa hồ sơ; nộp đơn nghỉ phép; đăng ký đi công tác; xem văn bản đến; theo dõi nhiệm vụ; điểm danh họp/báo vắng. \\
    \hline
    \textbf{Lãnh đạo Đơn vị} (Trưởng/Phó Khoa, Phòng) & \texttt{dv:nghi\_phep:read/write}, \texttt{dv:di\_cong\_tac:read/write}, \texttt{eofficeVanBanDen:manage}, \texttt{iofficeMission:manage} & Thẩm định \& phê duyệt đơn nghỉ phép đơn vị; duyệt danh sách đi công tác; phân phối chỉ đạo văn bản đến; giám sát tiến độ nhiệm vụ đơn vị; tạo lịch họp đơn vị. \\
    \hline
    \textbf{Ban Giám hiệu} (Hiệu trưởng, Phó HT) & \texttt{tcns:quy\_trinh:manage}, \texttt{eofficeVanBanDi:read}, \texttt{scheduleGeneral:manage} & Xem xét văn bản đi cấp trường; chỉ đạo văn bản đến quan trọng; theo dõi nhiệm vụ trọng tâm trường; duyệt lịch tuần trường. \\
    \hline
    \textbf{Chuyên viên TCCB} (Mã đơn vị 94) & \texttt{tcns:ly\_lich:manage}, \texttt{tcns:request\_ly\_lich:read}, \texttt{tcns:nghi\_phep:write} & Thẩm định Diff sửa lý lịch cán bộ; duyệt đơn nghỉ phép bước cuối \& cấp số quyết định; quản lý quỹ phép năm. \\
    \hline
    \textbf{Văn thư Trường/ĐV} & \texttt{eofficeVanBanDen:write}, \texttt{scheduleGeneral:write} & Nhập metadata văn bản đến; theo dõi luân chuyển văn bản đi; tổng hợp lịch tuần trường. \\
    \hline
    \end{tabularx}
    \caption{\textit{Ma trận Phân quyền người dùng theo vai trò (RBAC Matrix)}}
    \label{tab:rbac_matrix}
\end{table}

\section{Yêu cầu Chức năng (Functional Requirements)}
\label{sec:yeucau_chucnang}

\subsection{Phân hệ Xác thực và Quản lý Phiên (Auth Module)}
\begin{itemize}
    \item Đăng nhập qua HCMUT CAS SSO và xác thực Mật khẩu nội bộ.
    \item Quản lý Token mã hóa an toàn trong Keychain/Keystore.
    \item Tự động làm mới phiên làm việc ngầm (Silent Refresh) khi gặp mã lỗi 401 Unauthorized.
\end{itemize}

\subsection{Phân hệ Quản trị Nhân sự (HRM Module)}
\begin{itemize}
    \item \textbf{Hồ sơ Lý lịch:} Tra cứu 11 nhóm thông tin lý lịch chi tiết; gửi đề xuất chỉnh sửa kèm ảnh minh chứng; giao diện so sánh sai khác (Diff Viewer) dành cho người thẩm định.
    \item \textbf{Đăng ký Nghỉ phép:} Form Wizard 4 bước kết hợp thuật toán kiểm tra ràng buộc thời gian thực 4 tầng (tính số ngày làm việc, kiểm tra trùng lịch, mốc nộp trễ 3 ngày/15 ngày, kiểm tra số dư quỹ phép năm).
    \item \textbf{Đăng ký Đi công tác:} Nhập mục đích, địa điểm, thời gian, dự toán kinh phí, chọn thành viên đoàn công tác và đính kèm thư mời/kế hoạch PDF.
    \item \textbf{Phê duyệt Đa cấp \& Hàng loạt:} Hiển thị danh sách chờ duyệt, timeline quy trình luân chuyển; hỗ trợ duyệt nhiều đơn cùng lúc qua thanh tác vụ nhóm (\texttt{AppBatchActionBar}).
\end{itemize}

\subsection{Phân hệ Văn phòng số (iOffice Module)}
\begin{itemize}
    \item Tra cứu, tìm kiếm và phân loại danh mục văn bản đến, văn bản đi.
    \item Xem trực tiếp tệp PDF văn bản đính kèm qua trình xem tích hợp.
    \item Lãnh đạo đơn vị thực hiện phân phối, giao việc và bút phê chỉ đạo văn bản đến.
\end{itemize}

\subsection{Phân hệ Quản lý Nhiệm vụ \& Công việc (Tasks / Missions Module)}
\begin{itemize}
    \item Phân loại 3 nhóm nhiệm vụ (Trọng tâm, Được giao, Chung) kết hợp lọc theo 5 trạng thái (Cần xử lý, Nháp, Đang làm, Đã xong, Tất cả).
    \item Xem chi tiết nhiệm vụ với 4 tab (Tổng quan \% tiến độ, Cây đầu việc \texttt{outlined-tree}, Báo cáo tiến độ đợt, Văn bản liên kết).
    \item Modal xem chi tiết công việc phân công và gửi báo cáo tiến độ đính kèm minh chứng.
\end{itemize}

\subsection{Phân hệ Lịch công tác và Điểm danh Cuộc họp (Schedule Module)}
\begin{itemize}
    \item Tra cứu lịch tuần trường và lịch tuần đơn vị.
    \item Kiểm tra ràng buộc khung giờ mở điểm danh (từ trước giờ họp 1 giờ đến hết ngày diễn ra cuộc họp).
    \item Điểm danh có mặt hoặc báo vắng có lý do (1--200 ký tự); đồng bộ dữ liệu điểm danh thời gian thực qua WebSocket.
\end{itemize}

\subsection{Phân hệ Trung tâm Thông báo Đẩy (Notification Module)}
\begin{itemize}
    \item Tiếp nhận thông báo đẩy từ Firebase FCM; cập nhật số huy hiệu Badge trên icon ứng dụng.
    \item Phân tích 4 trường Metadata có cấu trúc (\texttt{source}, \texttt{entityType}, \texttt{entityId}, \texttt{isApproval}) để thực hiện điều hướng sâu (Deep Linking) trực tiếp vào màn hình chi tiết.
\end{itemize}

\section{Yêu cầu Phi chức năng (Non-Functional Requirements)}
\label{sec:yeucau_phichucnang}

\begin{itemize}
    \item \textbf{Hiệu năng và Độ mượt mà (Performance):} Tốc độ khung hình duy trì ổn định 58--60 FPS khi cuộn danh sách lớn; thời gian khởi động nguội $< 1.5$ giây; thời gian phản hồi API trung bình $< 150$ms.
    \item \textbf{Độ trễ Thông báo Đẩy (Latency):} Độ trễ đầu-cuối từ khi sự kiện phát sinh trên Backend đến khi thiết bị nhận Push Notification $< 2.0$ giây.
    \item \textbf{Bảo mật (Security):} HTTPS/TLS toàn diện; mã hóa Token trong Keychain/Keystore; kiểm tra phân quyền RBAC đa tầng ở cả Client và Server.
    \item \textbf{Khả năng Chịu lỗi (Fault-Tolerance):} Bắt lỗi mất kết nối mạng và Timeout nhẹ nhàng qua Snackbar, cung cấp nút Thử lại (Retry), tuyệt đối không làm dừng đột ngột ứng dụng (Crash).
\end{itemize}
